\documentclass[10pt,a4paper]{article}

\usepackage[a4paper,margin=0.7in]{geometry}
\usepackage[T1]{fontenc}
\usepackage[utf8]{inputenc}
\usepackage{lmodern}
\usepackage[hidelinks]{hyperref}
\usepackage{enumitem}
\usepackage{titlesec}
\usepackage{parskip}
\usepackage{fontawesome5}
\usepackage{xcolor}

\pagestyle{empty}
\setlength{\parindent}{0pt}
\setlist[itemize]{leftmargin=1.2em,itemsep=0.2em,topsep=0.15em,parsep=0em,partopsep=0em}

\definecolor{linkblue}{HTML}{1F4E79}
\hypersetup{
    colorlinks=true,
    urlcolor=linkblue
}

\titleformat{\section}{\large\bfseries}{}{0em}{}[\titlerule]
\titlespacing*{\section}{0pt}{0.8em}{0.4em}

\newcommand{\resumeheading}[4]{%
  \textbf{#1} \hfill #2\\[-0.15em]
  \textit{#3} \hfill \textit{#4}\\[-1.5em]
}

\begin{document}

\begin{center}
	{\LARGE \textbf{Shreyas Kalvankar}}\\[0.3em]
	% \faMapMarker*~Delft, Netherlands
	% \quad \textbar \quad
	\faPhone~+31 616243845
	\quad \textbar \quad
	\faEnvelope~\href{mailto:shreyaskalvankar@gmail.com}{shreyaskalvankar@gmail.com}\\[0.25em]

	\faGlobe~\href{https://obi-wan-shinobi.github.io}{obi-wan-shinobi.github.io}
	\quad \textbar \quad
	\faLinkedin~\href{https://linkedin.com/in/shreyas-kalvankar}{linkedin.com/in/shreyas-kalvankar}
	\quad \textbar \quad
	\faGithub~\href{https://github.com/obi-wan-shinobi}{github.com/obi-wan-shinobi}
\end{center}
% \section*{Summary}
% Machine Learning Research Engineer and Software Engineer with experience building compact neural networks, PyTorch training pipelines, and production-oriented ML systems for edge and resource-constrained environments. Worked on computer vision, distributed training, and data pipelines across research and applied settings.

\section*{Education}

\resumeheading{Technische Universiteit Delft}{September 2024 -- July 2026}{MSc Computer Science -- Machine Learning Applications Engineering \& Algorithmics}{Delft, Netherlands}
\begin{itemize}
	\item Master’s thesis: Neural network training dynamics and spectral bias

	      Investigated spectral bias and training dynamics of overparameterized neural networks, combining NTK theory, Fourier/operator analysis, and numerical experiments to study how finite width and sampling affect feature learning.
\end{itemize}

\resumeheading{Savitribai Phule Pune University}{August 2017 -- May 2021}{Bachelor of Engineering in Computer Engineering}{Pune, India}
\begin{itemize}
	\item Bachelor's thesis: Astronomical Image Colorization and Super-resolution using GANs
\end{itemize}
\section*{Experience}

\resumeheading{Technische Universiteit Delft}{September 2025 -- Present}{Teaching Assistant}{Delft, Netherlands}
\begin{itemize}
	\item Supported MSc courses in Machine and Deep Learning, Statistical Learning, Security and Cryptography, and Constraint Solving.
	\item Led tutorials, created and graded assignments and exams, and assisted students with debugging implementations.
\end{itemize}

\resumeheading{Dalton Maag Ltd.}{November 2021 -- February 2024}{Software Developer}{London, United Kingdom}
\begin{itemize}
	\item Automated CJK (Chinese-Japanese-Korean) glyph generation using a Python genetic-algorithm POC, demonstrating reduction of production time from \textbf{months to minutes}.
	\item Developed core features for PriceBot (TypeScript, Ruby on Rails) and productionized it, cutting quote generation time from \textbf{hours to seconds}.
	\item Designed a containerized GitLab CI/CD pricing health monitor with daily metric scraping and threshold alerting via emails that prevented silent regressions and potential six-figure revenue losses.
	\item Engineered a Devanagari complexity model and overhauled Arabic, Greek, and Cyrillic models using TypeScript and Ruby on Rails/PostgreSQL, scaling pricing coverage from \textbf{about 300 to about 2,500 glyphs}.
	\item Modeled effort-scaling heuristics and non-linear design constraints as a dependency graph to automate task scheduling, replacing a fully manual estimation process and reducing planning time from \textbf{hours to minutes}.
	\item Contributed to DSedit, an internal ElectronJS, VueJS, and Django platform, abstracting complex Unified Font Object (UFO) and Designspace specifications into a reliable GUI to eliminate manual XML editing and compilation errors.
\end{itemize}

\resumeheading{Relfor Labs Pvt. Ltd.}{September 2022 -- Present}{Machine Learning Research Scientist (Consulting, Part-time)}{Pune, India}
\begin{itemize}
	\item Designed compact, decomposable CNN architectures in PyTorch for edge-oriented deployments, enabling low-latency audio transition detection by caching hierarchical time-series embeddings during rolling-window inference.
	\item Built PyTorch Lightning and Distributed Data Parallel training pipelines on NVIDIA DGX A100, accelerating experimentation by \textbf{about 50\%}.
	\item Accelerated terabyte-scale (13TB) data preprocessing by \textbf{8x} using Numba JIT compilation and Pandas to optimize CPU utilization.
	\item Reduced pipeline storage footprint by \textbf{30\%} by migrating UTF-8 data to Apache Arrow, implementing byte-addressable binary storage for rapid memory loading.
\end{itemize}

\resumeheading{Relfor Labs Pvt. Ltd.}{August 2021 -- November 2021}{Machine Learning Engineer}{Pune, India}
\begin{itemize}
	\item Designed CNN architectures for audio classification using mel-spectrograms in PyTorch, steering experimentation to surpass standard baselines, achieving \textbf{about 98.6\%} accuracy and \textbf{about 0.98} F1-score.
\end{itemize}


\section*{Technical Skills}
\textbf{Languages:} Python, JavaScript, TypeScript, C, C++, Rust, Go\\
\textbf{Machine Learning:} PyTorch, TensorFlow, Keras, JAX, Optuna\\
\textbf{Backend and Web:} Ruby on Rails, React, Vue, Django, HTML, CSS\\
\textbf{Data and Databases:} PostgreSQL, MongoDB, Pandas, Apache Arrow\\
\textbf{DevOps and Infrastructure:} Docker, CI/CD, container orchestration, observability\\
\textbf{Cloud and Compute:} Google Cloud Platform, NVIDIA DGX A100, Distributed Data Parallel

\section*{Projects}

\resumeheading
{BingeBuddy}
{}
{Course Project --- \href{https://github.com/kanta8864/binge-buddy}{github.com/kanta8864/binge-buddy}}
{}
\begin{itemize}
	\item Built BingeBuddy, a task-oriented conversational recommendation system for personalized movie and TV series discovery.
	\item Designed the agent memory architecture to track and recall relevant information from past conversations, using LLM-based workflows to decide what should be stored in long-term memory, what to extract, and how to categorize and merge new information with existing memories.
	\item Developed the Django backend and used LangChain to orchestrate the agentic workflows for recommendation and memory management.
	\item Used MongoDB as the long-term memory store and ReactJS for the frontend, enabling more personalized multi-turn recommendations based on prior interactions.
\end{itemize}

\resumeheading
{REMLA Project}
{}
{Course Project --- \href{https://github.com/remla25-team15/operation}{github.com/remla25-team15/operation}}
{}
\begin{itemize}
	\item Built and deployed a production-oriented machine learning application from scratch with a focus on reproducibility and scalability.
	\item Used Vagrant, Ansible, Kubernetes, Helm, Istio, DVC, and Git-based CI/CD to manage infrastructure, versioning, and deployment workflows.
	\item Developed experience with release engineering practices for machine learning systems in a team-based software setting.
\end{itemize}

\resumeheading
{Galaxy Morphology Classification}
{}
{Research Project --- \href{https://github.com/obi-wan-shinobi/GalaxyEfficientNets}{github.com/obi-wan-shinobi/GalaxyEfficientNets}}
{}
\begin{itemize}
	\item A galaxy morphology classification, based on Kaggle Galaxy Zoo 2 competition, implementing the EfficientNet architectures in Tensorflow.
	\item Built TensorFlow CNNs for galaxy morphology prediction, achieving a top-3 Kaggle ranking for vote-fraction regression with RMSE 0.07765.
	\item Developed a 7-class galaxy morphology classifier with 93.7\% accuracy and 0.8857 F1-score.
\end{itemize}

\resumeheading
{EinsteinPy}
{}
{Open Source --- \href{https://github.com/einsteinpy/einsteinpy}{github.com/einsteinpy/einsteinpy}}
{}
\begin{itemize}
	\item Contributed core general relativity metrics and numerical fixes to a community Python library used for scientific computing and education.
	\item Added the Reissner–Nordström metric, a static solution to the Einstein–Maxwell field equations.
	\item Fixed issues in the Kerr–Newman and Kerr metric classes.
	\item Implemented calculations for the event horizon and ergosphere of a Kerr–Newman black hole.
\end{itemize}

\section*{Publications \& Pre-prints}
\begin{itemize}
	\item \textbf{\href{https://dl.gi.de/items/0ff89330-3adf-49b3-8071-53eb7f176488}{Astronomical Image Colorization and Up-scaling with Conditional Generative Adversarial Networks}}, INFORMATIK 2022
	\item \textbf{\href{https://arxiv.org/abs/2005.11288}{EinsteinPy: A Community Python Package for General Relativity}}, arXiv:2005.11288
	\item \textbf{\href{https://arxiv.org/abs/2008.13611}{Galaxy Morphology Classification using EfficientNet Architectures}}, arXiv:2008.13611
\end{itemize}

\section*{Languages}
English (Fluent, IELTS 8.5, CEFR C2), Dutch (Elementary), Japanese (Elementary)

\end{document}
